\section{The CDM Product Model---A Developer's Guide}
\label{sec:appF}
\label{app:cdm-walkthrough}
%==============================================================================

The ISDA Common Domain Model defines a layered object hierarchy for financial products. Each layer maps to one ledger construct, and the map crystallises a single result: the ledger unit is the CDM \texttt{Trade}, not the product. This appendix establishes that result through the six layers, then exercises three boundaries the basic layers do not reach---payout composition, tokenised custody, and date resolution---with one worked instrument per result. Throughout, ``CDM'' denotes Common Domain Model v6.0.0 (FINOS-governed); field names follow the Rosetta/Rune schema.

%----------------------------------------------------------------------
\subsection{The Six-Layer Hierarchy}
\label{sec:cdm-layers}
%----------------------------------------------------------------------

The layers narrow from shared to unique. Five are reusable infrastructure; only the Payout is product-specific, and only the Trade is a unit.

\paragraph{Layer 1 --- Observable.} The market quantity a contract references (an equity index, a rate). It is identical across product type, direction, and venue: an option and a future on the same index share one Observable. It is a market-data input to pricing and lifecycle observation---not a unit, never held or transferred.

\paragraph{Layer 2 --- Payout.} A choice type, and the only truly product-specific layer. An option is one \texttt{OptionPayout}; a future is a \texttt{PerformancePayout} paired with a \texttt{FixedPricePayout}, mirroring an equity swap. CDM does not distinguish listed from OTC here---the economic logic of a listed and an OTC option is the same payout with different strike, expiry, and settlement currency. A put flips \texttt{optionType}; a short reverses direction (\texttt{buyerSeller} for options, \texttt{payerReceiver} for swap-like legs); nothing else changes. The Payout is the move-generating logic of the smart contract: it defines what transfers occur under what conditions, and carries \emph{no} counterparty identity and \emph{no} collateral reference.

\paragraph{Layer 3 --- EconomicTerms.} {\emergencystretch=3em Wraps the Payout with \texttt{effective\allowbreak Date}, \texttt{termination\allowbreak Date}, and calculation-agent metadata. It drives the lifecycle state machine: the effective date triggers ACTIVE, the termination date triggers expiry, exercise, or settlement observation. Its \texttt{collateral} field is product-intrinsic (repo haircuts, securities-lending collateral)---\allowbreak not the variation-margin CSA, which attaches at Trade.\par}

\paragraph{Layer 4 --- NonTransferableProduct.} Wraps EconomicTerms with identifiers and ISDA taxonomy. It is the contract \emph{template}---the code, not the instance---shared by every trade with identical terms. CDM offers two product types: \texttt{Transferable\allowbreak Product} (extends \texttt{Asset}; for instruments whose economics follow from ownership---equities, plain bonds) and \texttt{Non\allowbreak Transferable\allowbreak Product} (wraps EconomicTerms; for contracts whose payout must be explicitly modelled). Listed options and futures are \texttt{Non\allowbreak Transferable\allowbreak Product}: exchange-traded and fungible, but contracts with explicit payout logic, not assets transferred by ownership. Their exchange fungibility is captured by a shared template---all trades in one contract specification reference the same NonTransferableProduct.

\paragraph{Layer 5 --- TradableProduct.} Binds a template to a \texttt{TradeLot} (agreed price and quantity) and counterparty roles. It is the instance before collateral attachment: counterparties and economics, no collateral terms. A product quoted at many prices to many parties is many TradableProducts; confirmation makes one a Trade.

\paragraph{Layer 6 --- Trade.} Extends TradableProduct with \texttt{tradeIdentifier} (UTI), \texttt{tradeDate}, resolved parties (LEIs), \texttt{executionDetails}, and \texttt{collateral}. The instruments share everything up to here and diverge completely at Trade: counterparty (named bilateral vs.\ CCP), collateral source (bilateral CSA vs.\ CCP rulebook), and venue.

\begin{center}\small
\begin{tabular}{lll}
\toprule
\textbf{CDM layer} & \textbf{Ledger construct} & \textbf{Varies across instruments?} \\
\midrule
Observable & Market-data input & No \\
Payout & Move-generating logic & \textbf{Yes---the only product-specific layer} \\
EconomicTerms & Contract definition & No (structural wrapper) \\
NonTransferableProduct & Contract template & Shared if terms identical \\
TradableProduct & Priced instance & Trade-specific \\
Trade & \textbf{Unit} & Always \\
\bottomrule
\end{tabular}
\end{center}

%----------------------------------------------------------------------
\subsection{Unit Identity: Trade Is the Unit}
\label{sec:cdm-trade}
%----------------------------------------------------------------------

\begin{principle}[Trade $=$ Unit]
The ledger unit is the full CDM Trade object, including its \texttt{Collateral} field. Product identity is not unit identity.
\end{principle}

The unit key is the Trade identifier (UTI for OTC; contract specification for listed). The forgetful map $U:\mathbf{Trade}\to\mathbf{NonTransferableProduct}$ that strips counterparty and collateral is many-to-one---meaning many distinct trades collapse to one product. The ledger operates on the source $\mathbf{Trade}$, never on the quotient, and the Unit Store (Section~\ref{sec:unit-store}) enforces injectivity at registration: distinct Trades map to distinct units.

\paragraph{The dual-CSA test.} Two OTC trades reference the same template---a European call on the same index at the same strike---but one carries a USD CSA (SOFR discounting, New York law) and the other a EUR CSA (ESTR discounting, English law). Because \texttt{Collateral} differs, the Trade objects differ, and the units differ. Two otherwise-identical options collateralised in different currencies are \textbf{different units}: different discount curves, margin flows, and settlement currencies.

\paragraph{Listed and cleared.} For listed instruments the CCP novates itself in as sole counterparty, so all positions in one contract specification are fungible regardless of original counterparty. Listed-vs-OTC is therefore a Trade-level distinction---counterparty becomes the CCP, collateral becomes the CCP rulebook, venue becomes the exchange---not a product-level one. For futures, daily settlement is itself the PnL realisation rather than collateral against it; the cash mechanics are in the futures lifecycle (Section~\ref{sec:sec08}).

\paragraph{CCP margin is a CDM gap.} CDM v6.0.0 has no type for CCP margin rules distinct from bilateral collateral provisions; the \texttt{Collateral} type on Trade was built for bilateral CSAs. For cleared trades the margin framework is an external reference to the CCP rulebook. The two-contract design (payoff contract plus margin contract) handles both regimes; CDM's representation of the cleared case is the less mature.

%----------------------------------------------------------------------
\subsection{Payout Composition and the TransferableProduct Boundary}
\label{sec:cdm-structured-note}
%----------------------------------------------------------------------

A structured equity note tests payout composition and the product-type boundary at once. BankCo issues a 1Y zero-coupon note on an equity index, EUR\,1{,}000{,}000 notional, to FundCo. Redemption is
\[
R \;=\; N \times \min\!\Bigl(1,\ \tfrac{S_T}{S_0}\Bigr),
\]
par if the index is flat or up, full downside if it falls. Economically the note is a zero-coupon bond \emph{plus} a short put struck at par. It has an ISIN, lists on an exchange, and settles through a CSD: a transferable security in law and market infrastructure.

\paragraph{Result 1: economic decomposition is not CDM decomposition.} The correct representation is a single \texttt{PerformancePayout} with a return capped at 0\%---not a \texttt{FixedPricePayout} (par bond) plus an \texttt{OptionPayout} (short put). The decomposition is economically correct yet wrong as a CDM model for three reasons: the embedded put has no independent exercise right (the redemption formula applies automatically, so there is no \texttt{OptionPayout}); the note has one contractual cash flow at maturity, not two; and CDM qualification recognises a capped \texttt{PerformancePayout} on an equity index as a known structure, but not a premium-less \texttt{OptionPayout}. Risk systems decompose into bond + put for Greeks; CDM models the contractual cash flow.

\paragraph{Result 2: \texttt{Transferable\allowbreak Product} is narrower than ``transferable.''} The note is a \texttt{Non\allowbreak Transferable\allowbreak Product} despite being a transferable security. The criterion is not legal transferability but modelling complexity: \texttt{Transferable\allowbreak Product} is for instruments whose cash flows follow mechanically from ownership (a plain bond's coupon and par); any instrument whose cash flows depend on an observable---requiring an explicit Payout tree---is \texttt{Non\allowbreak Transferable\allowbreak Product}. The ISIN lives in the \texttt{Non\allowbreak Transferable\allowbreak Product} \texttt{identifier} field; \texttt{Transferable\allowbreak Product} is not required to carry it. This parallels the IFRS~9 SPPI test: the note fails SPPI because of the equity-linked redemption and is measured at fair value, exactly the signal---non-trivial payout---that forces \texttt{Non\allowbreak Transferable\allowbreak Product} and trading-book scope (Section~\ref{sec:intro}).

\paragraph{Result 3: one unit, three settlement paradigms.} The note is one security, one ISIN, one Trade, one ledger unit; the embedded put is not separately tradable, exercisable, or settlable. Issuance and redemption are two-move transactions that conserve per unit. In the index-down-20\% scenario redemption is EUR\,800{,}000:
\[
\text{Move}(\text{BankCo}\to\text{FundCo},\ \text{EUR},\ 800{,}000),\qquad
\text{Move}(\text{FundCo}\to\text{BankCo},\ \text{note},\ 1).
\]
The ``missing'' EUR\,200{,}000 is the put payoff, already recognised through mark-to-market: before maturity FundCo holds one note valued at EUR\,800{,}000, after maturity EUR\,800{,}000 cash, and the value jump at redemption is zero---lifecycle value invariance (Section~\ref{sec:intro}) holds. The note also exposes the settlement spectrum: bilateral CSA (well modelled), CCP margin (partially modelled, Section~\ref{sec:cdm-trade}), and CSD settlement (essentially unmodelled---CDM v6.0.0 has no \texttt{CSDSettlement} type, so settlement through Euroclear or Clearstream is carried in generic \texttt{executionDetails}). CDM support decreases left to right, reflecting its OTC-derivative origin.

%----------------------------------------------------------------------
\subsection{Tokenised Securities: Double-Counting and the Flat Custodian}
\label{sec:cdm-tokenized}
%----------------------------------------------------------------------

Where the structured note tested product-type, a tokenised equity tests custody. In the mirror/backed model a regulated custodian holds real shares in segregation and issues on-chain tokens 1:1; each token is a beneficial claim on one share, held against the custodian. The underlying exposure, dividends, and corporate actions are unchanged; the identifier scheme (contract address, chain ID, token standard), custodial layer, settlement (on-chain atomic vs.\ T+1), and fractional divisibility are new.

\paragraph{Two units, not one.} The share (NVDA) and the token (NVDA\_TOKEN) are distinct units in the Unit Store: different settlement mechanisms, different identifiers, and different risk profiles---the token carries custodian credit risk the share does not. Treating them as one unit would erase the distinction between a holder of real shares and a holder of tokenised claims.

\paragraph{The double-counting problem and its resolution.} Distinct units invite a na\"ive aggregation of NVDA exposure that double-counts: 1{,}000{,}000 real shares in custody plus 1{,}000{,}000 tokens held by investors sums to 2{,}000{,}000, though only 1{,}000{,}000 shares of exposure exist.

\begin{center}\small
\begin{tabular}{@{}lrrr@{}}
\toprule
\textbf{Entity} & $\mathbf{own(NVDA)}$ & $\mathbf{own(NVDA\_TOKEN)}$ & \textbf{Net NVDA exposure} \\
\midrule
Custodian & $+1{,}000{,}000$ & $-1{,}000{,}000$ & $0$ \\
Holders (all) & $0$ & $+1{,}000{,}000$ & $+1{,}000{,}000$ \\
\midrule
\textbf{System} & $+1{,}000{,}000$ & $0$ & $+1{,}000{,}000$ \\
\bottomrule
\end{tabular}
\end{center}

\begin{principle}[The custodian is flat]
A backing custodian's net economic exposure is zero by conservation, not by choice. It is long the underlying and short the tokens it issued in equal measure: $own(\text{NVDA})=+1{,}000{,}000$ and $own(\text{NVDA\_TOKEN})=-1{,}000{,}000$, so a unit price move changes asset and liability identically.
\end{principle}

Economic exposure is the read-time projection $own(\text{NVDA})+own(\text{NVDA\_TOKEN})$ for any entity: zero for the custodian, $+100$ for a holder of 100 tokens. The token balance conserves to zero across issuer and holders; the share balance conserves against the virtual issuance wallet. If the custodian issued tokens without holding shares, the shortfall would surface as a conservation violation---the resolution is a consequence of system closure, not a modelling convention. The structure is the securities-lending custody problem (Section~\ref{sec:gpm}): possession is not ownership. SBL distinguishes them by the position vector's \texttt{own} and \texttt{onloan} coordinates; tokenisation distinguishes them by liability netting; the principle is identical.

\paragraph{Variants.} Two cases avoid the problem because no backing inventory exists: a \emph{synthetic} token is a delta-one derivative on the underlier (\texttt{Non\allowbreak Transferable\allowbreak Product} with a \texttt{PerformancePayout}, no custodian); a \emph{native-issuance} token \emph{is} the share on chain (\texttt{Transferable\allowbreak Product} with blockchain-native identifiers). The mirror token is the first \texttt{Transferable\allowbreak Product} in the walkthrough and the only one with a non-trivial custody layer.

\paragraph{Four CDM v6.0.0 gaps for tokenised securities.}
\begin{enumerate}[nosep]
\item \textbf{No blockchain identifier type.} \texttt{ProductIdentifier}/\texttt{ProductIdSourceEnum} omit contract address, chain ID, and token standard. Workaround: \texttt{Other} with a scheme convention. Significant---every token needs an ad-hoc convention.
\item \textbf{No custodial backing model.} No type relates a token to the underlying held by a custodian; backing ratio, custodian identity, and segregation are not representable. Workaround: trade-level \texttt{contractDetails}. Significant---the custodial relationship is the defining feature.
\item \textbf{\texttt{DigitalAsset} underspecified.} The branch exists but neither extends \texttt{Transferable\allowbreak Product} nor carries \texttt{EconomicTerms}; built for crypto-native assets, not tokenised securities. Moderate.
\item \textbf{No on-chain settlement type.} \texttt{SettlementTypeEnum} has no on-chain atomic value and no structure for chain, confirmation depth, or gas. Workaround: \texttt{Other}. Moderate.
\end{enumerate}
The gaps are incremental extensions, consistent with CDM's trajectory from bilateral OTC to cleared to tokenised; none requires an architectural change. On-chain delivery-versus-payment, when both legs settle on chain, is the strongest DvP---atomic within one transaction, eliminating settlement risk; mixed on/off-chain settlement degrades to conventional DvP risk, a settlement-layer concern (Section~\ref{sec:settlement}), not a ledger-layer one.

%----------------------------------------------------------------------
\subsection{Date Handling: Types, Resolution, and Gotchas}
\label{sec:cdm-dates}
%----------------------------------------------------------------------

A contract date is never a bare date: it carries adjustment rules, calendar references, and resolution logic that fix the operational date. CDM composes these from small precise types in the \texttt{base.datetime} namespace. The types feed every dated computation in this document---the IRS move schedule and day-count fractions (Section~\ref{sec:contracts}), coupon adjustment (Section~\ref{sec:temporal-scheduler}), and settlement-date adjustment (Section~\ref{sec:settlement}). A wrong calendar, convention, or day count is a valuation error, not a rounding error.

\subsubsection{The Type Taxonomy}
\label{sec:cdm-date-types}

\begin{center}\small
\begin{tabular}{@{}l >{\raggedright\arraybackslash}p{0.60\textwidth}@{}}
\toprule
\textbf{Type / enum} & \textbf{Role} \\
\midrule
\texttt{date} & Resolved ISO~8601 calendar date; the output, not the input. \\
\texttt{Period} & Multiplier + unit (D/W/M/Y); a zero-length period must be in days. \\
\texttt{Offset} & \texttt{Period} + \texttt{dayType} (Business/Calendar/Exchange\ldots); \texttt{dayType} present only when in days. \\
\texttt{BusinessDayConventionEnum} & FOLLOWING, MODFOLLOWING, PRECEDING, NEAREST, NONE\ldots \\
\texttt{BusinessCenters} & Holiday calendars (USNY, GBLO, JPTO\ldots); a date is good only if good in \emph{all} listed. \\
\texttt{BusinessDayAdjustments} & Convention + business centres: the reusable adjustment recipe. \\
\texttt{AdjustableDate} & \texttt{unadjustedDate} (contractual) + adjustments $\to$ \texttt{adjustedDate} (operational). \\
\texttt{RelativeDateOffset} & ``$n$ days/months from a reference date, adjusted by $C$ against $B$.'' \\
\texttt{AdjustableOrRelativeDate} & Required-choice union of the two above; used throughout the product model. \\
\texttt{CalculationPeriodFrequency} & Frequency + \texttt{rollConvention} (EOM, IMM, \_1..\_30); drives schedule generation. \\
\texttt{DayCountFractionEnum} & ACT/360, ACT/365F, 30/360, ACT/ACT (ISDA/ICMA)\ldots; each a dispatched \texttt{YearFraction}. \\
\bottomrule
\end{tabular}
\end{center}

\noindent MODFOLLOWING dominates OTC derivatives: it advances to the next business day unless that crosses a month boundary, then retreats, keeping the payment date aligned with its calculation period. Two day-count formulas recur: ACT/360 is $\text{actual days}/360$; 30/360 (Bond Basis) is
\[
\frac{360(Y_2-Y_1)+30(M_2-M_1)+(D_2-D_1)}{360},\qquad D_1=\min(d_1,30),\ D_2=\min(d_2,30)\text{ if }D_1>29.
\]

\subsubsection{The Resolution Chain}
\label{sec:cdm-date-resolution}

An \texttt{AdjustableOrRelativeDate} resolves deterministically in three steps; \texttt{ResolveAdjustableDate} encapsulates them.
\begin{enumerate}[nosep]
\item \textbf{Resolve the reference.} If relative, offset from the anchor (\texttt{date\allowbreak Relative\allowbreak To}): \texttt{Add\allowbreak Business\allowbreak Days} when \texttt{dayType = Business}, calendar arithmetic otherwise. Yields the unadjusted derived date.
\item \textbf{Apply business-day adjustment.} Test the unadjusted date against the centres; if not a good business day, shift by the convention.
\item \textbf{The adjusted date is final}---the operational date for the cash flow, exercise, or settlement.
\end{enumerate}
A 5Y USD SOFR swap rolling on the 3rd under MODFOLLOWING/USNY illustrates the chain: 2026-10-03 (Sat) $\to$ Mon Oct~5; 2027-10-03 (Sun) $\to$ FOLLOWING within month $\to$ Mon Oct~4; weekday roll dates pass unadjusted. The first fixed period (2026-04-03 to 2026-10-05, 30/360) gives $\delta=(30\cdot6+2)/360=182/360\approx0.5056$; the first floating period (2026-04-03 to 2026-07-03, ACT/360) gives $91/360\approx0.2528$. These adjusted dates and fractions are exactly the reset dates $t_k$ and $\delta_k$ the IRS contract consumes (Section~\ref{sec:contracts}).

\subsubsection{Gotchas}
\label{sec:cdm-date-gotchas}

Each is a bug class the type system prevents only when the types are used correctly.
\begin{enumerate}[nosep]
\item \textbf{Omitted business centres.} A convention with no centres adjusts against no calendar---every day is a business day and the convention is inert. Validate that any non-NONE convention has at least one centre.
\item \textbf{Calendar vs.\ business days.} ``T+2'' from a Friday is Sunday under Calendar, Tuesday under Business. \texttt{dayType} is load-bearing, not decoration.
\item \textbf{EOM vs.\ fixed-day roll.} \texttt{EOM} rolls to each month's last day (Feb 28/29, Mar 31); \texttt{\_28} always rolls to the 28th. Confusing them shifts payments by up to three days.
\item \textbf{Joint-holiday intersection.} Multiple centres take the \emph{intersection} of business days: good in New York but a London holiday is not good for USNY+GBLO. More dates are non-business under joint calendars than either alone.
\item \textbf{Stub periods.} An effective date off the roll cycle produces a short or long first period; \texttt{balanceOfFirstPeriod} places the stub. Ignoring it gives a wrong year fraction for that period.
\item \textbf{Day-count mismatch.} For a 182-day semi-annual period, ACT/360 gives $0.5056$ and 30/360 gives $0.5000$. On EUR/USD\,50{,}000{,}000 at 4\%, the difference is $0.005556\times0.04\times50{,}000{,}000\approx\$11{,}111$---a valuation error, not a rounding error.
\end{enumerate}
